% calculus.tex - The Lambda^Cap Calculus

\section{The \lambdaCap{} Calculus}
\label{sec:calculus}

\lambdaCap{} is a small extension of the simply-typed lambda calculus with three
capability constructs: %
capability access $\tmcvar{c}$, %
capability binding $\tmwith{c}{e_1}{e_2}$, %
and capability restriction $\tmallow{\varphi}{e}$. %

We assume the reader is already comfortable with standard formal-language
notions such as syntax, typing judgments, and big-step evaluation for the
simply-typed lambda calculus; Pierce's \emph{Types and Programming Languages} is
a standard reference \cite{pierce-tapl}. %

\autoref{fig:syntax} gives the reader-facing surface syntax. %
Function syntax uses brace notation, while function types place the designated
capability set above the arrow: $\tmabs{x}{T_1}{\varphi}{\varphi}{t}$ denotes a
function that receives the call-site capabilities in $\varphi$, and
$\tarr{T_1}{T_2}{\varphi}{\varphi}$ is the corresponding function type. %
The types of the capabilities named in $\varphi$ are defined in the global
capability declaration $\cpdecl$. %

\begin{figure}[t]
\centering
\begin{align*}
\text{Names} \quad & n \in \mathbb{N} \\[0.25em]
\text{Capability names} \quad & c, d \\[0.25em]
\text{Capabilities} \quad & \varphi ::= \effempty \mid \eff{c} \mid \varphi_1 \effunion \varphi_2 \\[0.25em]
\text{Types} \quad & T ::= \tnat \mid \tarr{T_1}{T_2}{\varphi}{\varphi} \\[0.25em]
\text{Terms} \quad & t ::= \tmconst{n} \mid x \mid \tmcvar{c} \mid
    \tmabs{x}{T_1}{\varphi}{\varphi}{t} \mid \tmapp{t_1}{t_2} \\
    & \quad\;\;\mid \tmwith{c}{e_1}{e_2} \mid \tmallow{\varphi}{e}
    \\[0.25em]
\text{Values} \quad & v ::= \vlconst{n} \mid
    \vlclosshort{x}{t}{\vlenv}{\cpenv}
\end{align*}
\caption{Syntax of \lambdaCap{}}
\label{fig:syntax}
\end{figure}

\subsection{Semantics}
\label{sec:semantics}

Evaluation uses a value environment $\vlenv$ for term variables and a capability
environment $\cpenv$ for capabilities. %
The evaluation judgment
\[
\eval{t}{\vlenv}{\cpenv}{v}
\]
means that $t$ evaluates to $v$ under those environments. %

The rules are shown in \autoref{fig:eval-rules}. %
Capability lookup reads from $\cpenv$. %
Capability binding evaluates a body under an updated capability environment. %
Restriction runs a term in a restricted capability environment. %
The main rule is closure creation: evaluating
$\tmabs{x}{T}{\varphi}{\varphi}{t}$ stores $\cpdiff{\cpenv}{\varphi}$, not the
full capability environment. %

\begin{figure}[t]
\centering
\begin{mathpar}
\inferrule*[Right=\rn{E-Const}]
  { }
  {\eval{\tmconst{n}}{\vlenv}{\cpenv}{\vlconst{n}}}

\inferrule*[Right=\rn{E-Var}]
  {\vlenv(x) = v}
  {\eval{x}{\vlenv}{\cpenv}{v}}

\inferrule*[Right=\rn{E-CVar}]
  {\cplookup{\cpenv}{c} = v}
  {\eval{\tmcvar{c}}{\vlenv}{\cpenv}{v}}

\inferrule*[Right=\rn{E-Abs}]
  { }
  {\eval{\tmabs{x}{T}{\varphi}{\varphi}{t}}{\vlenv}{\cpenv}
        {\langle\tmabs{x}{T}{\varphi}{\varphi}{t},\; \vlenv,\; \cpdiff{\cpenv}{\varphi}\rangle}}

\inferrule*[Right=\rn{E-App}]
  {\eval{t_1}{\vlenv}{\cpenv}{\langle\tmabs{x}{T}{\varphi}{\varphi}{t},\; \vlenv_1,\; \cpenv_1\rangle} \\
   \eval{t_2}{\vlenv}{\cpenv}{v_2} \\
   \eval{t}{(\vlenv_1, x \mapsto v_2)}{(\cpmerge{\cpenv}{\cpenv_1})}{v_3}}
  {\eval{\tmapp{t_1}{t_2}}{\vlenv}{\cpenv}{v_3}}

\inferrule*[Right=\rn{E-With}]
  {\eval{e_1}{\vlenv}{\cpenv}{v_1} \\
   \eval{e_2}{\vlenv}{\cpupdate{\cpenv}{c}{v_1}}{v_2}}
  {\eval{\tmwith{c}{e_1}{e_2}}{\vlenv}{\cpenv}{v_2}}

\inferrule*[Right=\rn{E-Allow}]
  {\eval{e}{\vlenv}{\cprestrict{\cpenv}{\varphi}}{v}}
  {\eval{\tmallow{\varphi}{e}}{\vlenv}{\cpenv}{v}}
\end{mathpar}
\caption{Evaluation rules for \lambdaCap{}}
\label{fig:eval-rules}
\end{figure}

This distinction corresponds to two different use cases. %
Some capabilities are meant to come from the function's creation context and
therefore should be captured in the closure. %
Others are meant to come from the call site and therefore should be supplied by
the caller at invocation time. %

% typing.tex - Type System and Capability Tracking

\subsection{Type System and Capability Tracking}
\label{sec:typing}

\autoref{fig:typing-rules} presents the full system. %
The type system assigns each term both a type and a capability set. %
The type
describes the value produced; the capability set is a static upper bound on
the authority the term may use. %
The judgment
\[
\typ{\cpdecl}{\tpenv}{t}{T}{\varphi}
\]
means that term $t$ has type $T$ and may access capabilities in $\varphi$. %

\begin{figure}[t]
\centering
\begin{mathpar}
\inferrule*[Right=\rn{T-Const}]
  { }
  {\typ{\cpdecl}{\tpenv}{\tmconst{n}}{\tnat}{\effempty}}

\inferrule*[Right=\rn{T-Var}]
  {\tpenv(x) = T}
  {\typ{\cpdecl}{\tpenv}{x}{T}{\effempty}}

\inferrule*[Right=\rn{T-CVar}]
  {c \in \cpdecl \\
   \cpdecl(c) = T}
  {\typ{\cpdecl}{\tpenv}{\tmcvar{c}}{T}{\eff{c}}}

\\

\inferrule*[Right=\rn{T-Abs}]
  {\typ{\cpdecl}{\tpenv, x{:}T_1}{t}{T_2}{\varphi'}}
  {\typ{\cpdecl}{\tpenv}{\tmabs{x}{T_1}{\varphi}{\varphi}{t}}
       {\tarr{T_1}{T_2}{\varphi}{\varphi}}{\varphi' \effdiff \varphi}}


\inferrule*[Right=\rn{T-App}]
  {\typ{\cpdecl}{\tpenv}{t_1}{\tarr{T_1}{T_2}{\varphi}{\varphi}}{\varphi_1} \\
   \typ{\cpdecl}{\tpenv}{t_2}{T_1}{\varphi_2}}
  {\typ{\cpdecl}{\tpenv}{\tmapp{t_1}{t_2}}{T_2}
       {\varphi \effunion \varphi_1 \effunion \varphi_2}}

\inferrule*[Right=\rn{T-With}]
  {\cpdecl(c) = T_1 \\
   \typ{\cpdecl}{\tpenv}{e_1}{T_1}{\varphi_1} \\
   \typ{\cpdecl}{\tpenv}{e_2}{T_2}{\varphi_2}}
  {\typ{\cpdecl}{\tpenv}{\tmwith{c}{e_1}{e_2}}{T_2}
       {\varphi_1 \effunion (\varphi_2 \effdiff \eff{c})}}
\\

\inferrule*[Right=\rn{T-Allow}]
  {\typ{\cpdecl}{\tpenv}{e}{T}{\varphi'} \\
   \varphi' \effsubset \varphi}
  {\typ{\cpdecl}{\tpenv}{\tmallow{\varphi}{e}}{T}{\varphi'}}

\inferrule*[Right=\rn{T-Sub}]
  {\typ{\cpdecl}{\tpenv}{t}{T}{\varphi} \\
   c \in \cpdecl}
  {\typ{\cpdecl}{\tpenv}{t}{T}{\varphi \effunion \eff{c}}}
\end{mathpar}
\caption{Typing rules for \lambdaCap{}}
\label{fig:typing-rules}
\end{figure}

\rn{T-CVar} is the rule that actually uses a capability: accessing $\tmcvar{c}$
adds exactly $c$ to the term's capabilities. %
\rn{T-Abs} is the central rule. %
If the function body may use capabilities $\varphi'$, and the lambda marks
$\varphi$ as caller-provided, then the lambda expression itself requires only
$\varphi' \effdiff \varphi$ at its definition site. %
In other words, capabilities listed in $\varphi$ are expected from the call site
rather than from closure creation. %
The function type records the caller-provided capability set $\varphi$ above the
arrow. %

\rn{T-App} restores caller responsibility: applying a function of type
$\tarr{T_1}{T_2}{\varphi}{\varphi}$ contributes $\varphi$ to the application's
capabilities. %
This is the static counterpart of the runtime semantics, where the closure does
not capture the capabilities in $\varphi$ and the caller must supply them at
application time. %

\rn{T-With} binds a value to capability $c$ and masks that requirement in the
body, so uses of $c$ inside $e_2$ do not remain obligations of the outer
context. %
\rn{T-Allow} turns dynamic restriction into a static contract: if $e$ checks
with capability set $\varphi'$ and $\varphi' \subseteq \varphi$, then
restricting execution to $\varphi$ cannot remove a capability that the term
needs. %
\rn{T-Sub} is standard weakening. %


\subsection{Security Properties}
\label{sec:security}

At a high level, \lambdaCap{} guarantees that
\begin{itemize}[nosep]
\item capability annotations are trustworthy upper bounds on capability use, and
\item restricted computations cannot rely on capabilities outside the
  declared boundary.
\end{itemize}

If a term is typed with capabilities $\varphi$, then evaluation may access only
capabilities in $\varphi$ from the enclosing environment.

\textbf{Confinement}. %
Because capability provisioning and access is explicit and checked, well-typed
code is confined to the provided capabilities.
